\section{The negative quadratic-eight label on the original window}
\label{sec:coefficient-three-negative8-window}

\paragraph{The complete preceding union is retained.}
Keep the actual host, original coefficient completions, finite datum,
page and prescribed pair of AC-G183-window. In particular,
\[
 \begin{gathered}
 H_u=\mathbb Q(\zeta_{2^{h(u)+2}},\zeta_7,
       \sqrt{-q},\sqrt{-d},\omega,\sqrt[3]{269}),
       \qquad K=\mathbb Q(\omega),\\
 u^{2+8\varepsilon}\leq2^{h(u)}\leq2u^{2+8\varepsilon},
 \qquad0<\varepsilon<1/256,
 \qquad T=\lfloor u^{399}\rfloor.
 \end{gathered}\tag{N3W.1}
\]
Write $Z^{\rm two,inert}$ and $I_{399}^{\rm two,inert}$ for the
exact omitted set and remainder of (I2W.3)--(I2W.7).
Every symbol $N_T,C(h),E_3(h),P(h),B_T,A_T,R_T,S_T,J_-(T)$
retains its meaning there; in particular $N_T=T-1$.
The set $\mathscr E_h(O_3)$ and its count $E_3(h)$ use the
unchanged original residue embedding in (R3W.2)--(R3W.4).

\paragraph{One additional original finite label.}
For $h\geq2$ define
\[
 \begin{aligned}
 Z^{-8}_3(u)&=
       \{(3,\chi_{-8}|_{G_K},e):1\leq e<T\},\\
 Z^{\rm three,-8}(u)&=
                  Z^{\rm two,inert}(u)\sqcup Z^{-8}_3(u).
 \end{aligned}\tag{N3W.2}
\]
The new label has trivial conductor-seven, auxiliary quadratic
and Kummer cubic factors. Let $I_{399}^{\rm three,-8}$ be the
original sum (TC.1) with precisely this union omitted.
The weight is the original Tate weight; it is not a new finite
character at the coefficient prime.

\begin{theorem}[AC-G188-window: the negative-eight coefficient-three row]
\label{AC-G188-window}
Every triple in (N3W.2) is an actual original-coefficient strict zero,
and its displayed union is disjoint. For integers $h\geq2$ and $T\geq1$,
\[
 \begin{aligned}
 |Z^{\rm three,-8}\setminus Z^{\rm two,inert}|&=N_T,\\
 |Z^{\rm three,-8}_3|&=(2+E_3(h))N_T,\\
 2+E_3(h)&=
 \begin{cases}8,&h=2,\\2^{h+1}-2,&h\geq3.
 \end{cases}
 \end{aligned}\tag{N3W.3}
\]
All counts at every other coefficient prime are exactly those of
(I2W.4)--(I2W.5). The cumulative count is therefore
\[
 \begin{split}
 |Z^{\rm three,-8}|={}&
 (10+C(h)+E_3(h)+P(h)+5\cdot2^{h+1})N_T\\
 &-B_T+12A_T-R_T-S_T+J_-(T).
 \end{split}\tag{N3W.4}
\]
For $h\geq3$, all pure two-power finite labels at coefficient three
are covered except the two even primitive conductor-$32$ labels
failing the original-residue condition in (R3W.3).
Those remaining strict groups are not asserted nonzero.
The enlargement of the omitted set is proper whenever $T\geq2$;
strict decrease of the numerical remainder is not asserted.

Let $Q_u$ be the normalized positive derived-index cost of the
\emph{entire} union (N3W.2). Then
\[
 \begin{aligned}
 I_{399}&=I_{399}^{\rm three,-8}+Q_u,\\
 I_{399}^{\rm three,-8}&\leq I_{399}^{\rm two,inert}\leq I_{399},\\
 0\leq Q_u&\leq8D_u^2T^2\log(8TD_u)
      +\tfrac12D_u\mathcal R_uT^2+2D_ur_uc_0T\\
 &=O_\varepsilon(u^{802+16\varepsilon}(\log u)^2).
 \end{aligned}\tag{N3W.5}
\]
One complete (FL.4) majorant covers every old and new subfamily
together. These are nonnegative extended sums; no individual
positive derived index is claimed to vanish. With the exact
(TC.3) quantifiers and positive relative slack, a bound for
$I_{399}^{\rm three,-8}$ supplies the full original numerical
premise on the identical common window. Every original coefficient,
full local condition and actual source rank $R.t$ is retained.
\end{theorem}

\begin{proof}
\proofstep{1. The literal new strict theorem on the original completion.}
The host in (N3W.1) contains $\mathbb Q(\zeta_8)$, so the quadratic
character $\chi_{-8}$ is an actual finite label. Its values $\pm1$
already belong to the prescribed completion $O_3$ of
$\mathbb Z[\zeta_{3\cdot2^h}]$. With $E_3^{\rm fld}=\operatorname{Frac}(O_3)$,
the actual coefficient of (FC.2) is canonically
\[
 (O_3\otimes_{\mathbb Z_3}\mathbb Q_3/\mathbb Z_3)
                   (\chi_{-8})(-e)
      =(E_3^{\rm fld}/O_3)(\chi_{-8}c_3^{-e}).
 \tag{N3W.6}
\]
No replacement of this ring, lattice or embedding is made.
Apply the full strict theorem (G187.1)--(G187.2) with this
original $O_3$, the same positive $e$, and $\kappa=1$.
Its proof uses the two places over three of
$L=K(\sqrt{-2})$ to prove that one chosen decomposition group
surjects onto the maximal pro-three quotient. The finite
semidirect cocycle argument keeps the original uniformizer and
the actual local coboundary witness. Restriction back to $K$
is injective by the prime-to-three degree two.
Thus it supplies the original coefficient in (N3W.6) at every
positive weight, without a local-invariant divisibility hypothesis.

The strict conditions are full local zero above three and zero
in every away inertia group. Inflation from the original
$G_{K,\Sigma_3}$ in (FC.2) is injective on this same module.
Outside its defining ramification set an inflated cocycle kills
inertia already; at every place in the set the required local
condition is unchanged. Consequently the actual strict group
injects into the zero group of G187. All additional host primes
are retained. Every old triple is still supplied by G183.

\proofstep{2. Exact disjointness and the original-residue count.}
At coefficient three, the entire previous omitted set consists
of the trivial label and $\mathscr E_h(O_3)$ at every weight.
The latter contains $\chi_{-4},\chi_8$, all primitive conductor-$16$
characters, the specified conductor-$32$ characters, and all
primitive characters of conductor $2^c$ for $c\geq6$ in the host.
It does not contain $\chi_{-8}$.
Distinct rational two-power characters remain distinct upon
restriction to $G_K$: their ratio is unramified outside two and
cannot equal the nontrivial conductor-three character $\chi_{-3}$.
In particular $\chi_{-8}$ remains distinct from the trivial label
and from $\chi_8$ despite their equal conductor.
No other coefficient-three set was introduced by G183.
Different coefficient primes distinguish all the remaining old
triples, so the union in (N3W.2) is disjoint.

There are exactly $T-1$ positive integers below the positive
integer $T$. The new set therefore has cardinality $N_T$.
The equalities $E_3(2)=6$ and $E_3(h)=2^{h+1}-4$ for $h\geq3$
were proved with the actual residue embedding in (R3W.7).
Adding the single label proves (N3W.3); adding $N_T$ to (I2W.6)
proves (N3W.4). The two remaining conductor-$32$ labels are exactly
the same embedding-dependent exceptions as before.
When $h=2$, all eight pure two-power labels are now covered.
For $T\geq2$ at least the new triple of weight one is present,
which proves proper inclusion independently of any index value.

\proofstep{3. One normalized full-local error for the entire union.}
Use the original codifferent Tate-dual identification and
$b_\ell=\operatorname{rank}_{\mathbb Z_\ell}O_\ell$ from (FL.2).
The determinant inequality (EU.2), with its all-prime application,
gives for every omitted triple
\[
 \begin{aligned}
 \frac{\log\#k_{O_\ell}}{b_\ell}\max(0,i_{\ell,\eta,e})
 &\leq\frac{\log\#H^2(C_{\ell,\eta,e})}{b_\ell}\\
 &=\frac{\left|\log\#H^2(C_{\ell,\eta,e})
                 -\log\#\mathcal S_{\ell,\eta}(e)\right|}{b_\ell}.
 \end{aligned}
 \tag{N3W.7}
\]
The equality uses precisely the proved strict zero, so that the
strict group's order is one. It does not replace $H^2$ by zero.
At every weight, this whole set of omitted triples is one
nonnegative subsum of (FL.4), already summed over all coefficient
primes and all actual characters. Sum that bound once over
$1\leq e<T$, using $\sum e\leq T^2/2$, to obtain (N3W.5).
No second error is added for the new label or for old subfamilies.
The integral Fourier errors (FC.4)--(FC.5), in particular at
two and three, remain in the unchanged retained arithmetic interface.
There is no division by the order of the growing host.

The original bounds $D_u\leq288u^{2+8\varepsilon}$,
$r_u\leq12$ and
$\mathcal R_u=O_\varepsilon(u^{\varepsilon/5}(\log u)^2)$
give the stated exponent uniformly in the prescribed pair.
Partitioning a nonnegative sum proves the displayed additive
decomposition even when an unbounded remainder is infinite;
no subtraction of infinities occurs.

\proofstep{4. The identical window and unchanged remaining input.}
Fix a retained finite datum, page and $0<\varepsilon<1/256$.
Assume the exact (TC.3) premise for $I_{399}^{\rm three,-8}$:
there exist $J\geq1$, $X_0$ and $s\in(0,1)$, independent of
$X,q,d,u,\ell,\eta,e$, such that every prescribed pair at
$X\geq X_0$ has one $W$ with $X^{1/J}\leq W$ and $W^2\leq X$
on which the new remainder is at most
$(1-s)\exp(u^\varepsilon/2)$ for every $u\in[W,W^2]$.
The single uniform error (N3W.5) is
$o(\exp(u^\varepsilon/2))$. Choose $U_0\geq2$ making it at
most $(s/2)\exp(u^\varepsilon/2)$, and increase the threshold to
$\max(X_0,U_0^J)$. Keep exactly the supplied $W$.
Then $W\geq U_0$ and the full original premise holds throughout
that same interval with slack $s/2$.

Every old remainder premise implies the new one by nonnegativity.
Uniform little-$o$ premises and fixed finite simultaneous families
use this same proof with the maximum of their uniform constants;
separate windows are not intersected. The original page, pair,
host, coefficient completions, full local conditions and actual
$R.t$ are unchanged through the retained K35 implication.
The mixed theta/Selmer realization $\mathcal R_{\rm cub}$ is
still an input, with exactly its original coefficient and local
normalizations. No stronger AW.1 hypothesis is introduced.
\end{proof}

\paragraph{Boundary and retained dictionary.}
Only the one pure coefficient-three label $\chi_{-8}$ has been
added. All unsupported prime, character and weight combinations
remain in the remainder, including the quadratic-only modulo-seven
factor at coefficient two and every unproved mixed or auxiliary row.
The all-prime numerical bound, $\mathcal R_{\rm cub}$ and H04.4
remain open. The already audited SGA~1, Expos\'e~V, \S\S4,7 and
Stacks~03PK finite-cover/Artin--Tate dictionary is retained;
no new archive attribution, mixed motive or Lean execution is claimed.
